\chapter{Fairy Tale}
% Many years ago, in the faraway Valley of Siliconia, The Council of Data Pirates discovered that they needed quick means to analyze their records on which important trade routes to send their pirate ships to. 
% Over time consensus evolved to analyze the graphs by ingesting them from ETL pipelines into main memory on a large cluster.
% While fast, this approach limits the size of graphs that can be analyzed to the capacity of installed RAM available on the cluster and introduced the (often) unnecessary burden of managing a distributed system. 
% The alternative approach would be to store the data in a graph-native format persistently and stream it for analysis.
% Much effort has gone into creating systems that stream the graph from disk at speeds close to the available maximum I/O bandwidth, essentially re-inventing benefits that come from using a mature storage engine.

% We present \project{}, a graph database that allows storing the graph in formats that are amenable to high-ingest rates while also supporting analytics workloads and point-queries.
% The system itself is built as a thin Graph API layer on top of a key-value store, which allows us to experiment with different representations and enables us to support both vertex-centric and edge-centric style of writing algorithms.
% Our analysis shows which formats and algorithms are best suited for 
% We show that our \project{} is able to support high ingest rates and guarantees strong consistency and isolation, while remaining competitive with other out-of-core specialized graph processing systems.
% Moreover, since we have a unified analytics and transactional store that uses a read-friendly, we eliminate expensive pre-processing steps that are essential to other GPS.

% Villain: large data too big for RAM and preprocessing hard. Databases not used. GPS static. 
Once upon a time: people realized that their data can be represented as graphs and suddenly everybody wanted to analyze their large (and growing larger) graph datasets.
The systems community rose to the challenge and came up with many systems that let people analyze their large graph systems efficiently.
Many systems were proposed, each specialized in unique ways, which made them less generalizable for varying workloads.
However, most of these systems were designed to run as part of an ETL pipeline, ingesting raw data for each analysis.
This is an odd choice given the fact that the data being analyzed likely originated in a database and all the ingested data in its analysis-amenable respresentation is typically discarded when new data becomes available.

Over time the divide between the graph processing and the databases communities deepened, and there has been no convergence around a design that allows storage and analysis together.
Further driving this division is the fact that one-size-fits-all solutions do not work for graph processing systems any more than they do for databases.
Different graph processing problems exhibit different access patterns, and different datasets differ in key statistical properties (such as degree skew and diameter) making a universal choice of data representation for analysis and storage sub-optimal.
\PM{villain: dichotomy between graph processing and databases; one size does not fit all.}\\
We claim that a better solution combines a persistent representation with one amenable to efficient analytics.
Our system, \project{} is designed to leverage a mature key-value store to persist a graph in different formats, which allows a system designer to select one best suited to their design objectives.
We provide a thin API layer that provides basic graph storage and point-lookup primitives and iterator-based methods that provide efficient sequential access to the graph.
These APIs allow a developer to express many graph algorithms without having to worry about how the underlying storage layer is organized and without conforming to a restrictive vertex- or edge-centric style.
Our experiments comparing \project{} to popular specialized graph processing systems show that we remain performance competitive with them, while providing better data manageability.
We are faster than the graph databases we compared to, proving that \project{} represents a mid-point in the design space between databases and graph processing systems.